\section{Fazit}

In der vorliegenden Arbeit wurde gezeigt, wie die Maximierung einer aus
einem gleichschenkligen Dreieck und einem angelagerten Halbkreis
zusammengesetzten Fläche mithilfe eines chemischen Analogcomputermodells
implementiert werden kann. Die Simulationen in COPASI bestätigen, dass das Modell für verschiedene Eingabewerte die erwarteten Werte für den maximalen
Flächeninhalt und den normierten optimalen Winkel liefert.
Für den exemplarischen Eingabewert \(s=2\) lagen alle Ausgabewerte nach
etwa \(9.3\) Modellzeiteinheiten mit einer Genauigkeit von mindestens
\(99.9\,\%\) vor. 
Aufbauend auf diesem Modell könnten weitere geometrische
Optimierungsprobleme betrachtet werden, beispielsweise zusammengesetzte
Flächen mit anderen Kreissegmenten oder die chemische Berechnung
zusätzlicher geometrischer Kenngrößen.
